\scanpage{41}
Soit $K$ un 1-complexe\unsure{s|{}} combinatoire fini,
$(D_i)_{i\in J}$ \deleted{une} \deleted{\unsure{$f$}} ens. fini de sous-complexes
polygonaux, tels que toute arête de $K$ appartienne
à un $D_i$ et un seul, \deleted{et que $K$ n'ait pas
de sommet isolé i.e. $\forall$ sommet de $K$
appartient à au moins un des $D_i$.} On
suppose de plus que \deleted{si} $i\ne j$, $D_i$ et $D_j$
ont en commun un sommet et un
seul, et que $\forall$ sommet \inserted{de $K$} appartient à au
moins 2 des $D_i$. On dit alors que
$(K,(D_i)_{i\in J})$ est un arrangement \inserted{\deleted{\unsure{linéaire}}1-dimensionnel} de ps. droites,
\deleted{\unsure{$\forall$}} pourvu que l'intersection des $D_i$ soit vide.
i.e. que $K$ ne soit pas simplement le ``\unsure{point|joint}''
d'une famille de \deleted{polygones}\inserted{monogones (ou ``boucles'')}
rattachés par un sommet \inserted{\unsure{en|un}} commun. Il existe alors
tjrs au moins une partie $I\subset J$ de cardinal
3, tel que les $D_i$ $(i\in I)$ n'aient
pas de pt en commun. \deleted{$D_i$, \unsure{le}} \inserted{Pour}
% slanted
\margin{\uline{NB} L'ordre \unsure{$\nu_i$|ici} des polygônes $D_i$ \unsure{cf.} est \unsure{donc} \unsure{tjrs|tjs} $\geq 2$, pour un arrangem\textsuperscript{t} : aucun des $D_i$ n'est un monogône\ldots}%
\deleted{Pour} \unsure{tt|une} partie $I$ de $J$, soit $K_I$ le ss-complexe
de $K$ \inserted{qui ``topologiquement'' est} réunion des $D_i$ $(i\in I)$
% slanted
\margin{\uline{NB} $K_I$ sera regardé comme arrangem\textsuperscript{t} 1-dim de ps. droites, grâce aux $D_i$ $(i\in I)$, OK. si les $D_i$ ne sont pas \deleted{\unsure{con}} concourantes (il faut donc $\operatorname{card} I\geq 3$)}%
\inserted{mais où on prend comme sommets que les intersections des $D_i$ deux à deux.}
Dans le cas qui nous intéresse $(\operatorname{card} I=3,
\bigcap_{i\in I}D_i=\varnothing)$ \deleted{\unsure{on}} $K_I$ est un graphe
octogonal gauche, donc on peut parler
d'\uline{orientation} de ce graphe octogonal g.
% archive sheet number: 40
